\section{Overview}
\label{sec:overview}

Triton is a deterministic diagram compiler. It takes a small textual source and
produces a single, byte-stable SVG (optionally rasterised to PNG). Its surface
syntax is a \emph{superset of Mermaid}: existing Mermaid sources render
unchanged, and Triton adds diagram families Mermaid does not have.

\subsection{The central claim}

Every diagram kind --- whether a Mermaid flowchart or a Triton-native AVL tree
--- is a \emph{\texttt{DiagramModule}}: a unit that parses its own source into
its own intermediate representation (IR), lays that IR out, and emits a shared
\texttt{Scene}. A single renderer turns any \texttt{Scene} into SVG. There is no
``god IR'': each diagram owns the data model that fits it, and they meet only at
the \texttt{Scene} contract.

This contracts-first design is what makes the system both broad and consistent.
Adding a diagram kind means writing a parser and a layout function against fixed
contracts; it never means touching the renderer, the theme system, or other
diagrams.

\begin{figure}[htbp]
  \centering
\begin{triton}
poster "SQL Engine — Anatomy"
    columns 2

    cell q "Query plan (EXPLAIN)" :: plan
        plan
            Hash Join {rows: 980}
                Seq Scan orders {rows: 10000}
                Hash
                    Index Scan customers {idx: idx_cust}
    end

    cell idx "B+tree index" :: btree
        btree order 3 insert 10 20 30 40 50 60 25 55
    end

    cell pg "Slotted page" :: page
        page
            slots 3
            tuples (10,Ann) (40,Bob) (50,Cy)
    end

    cell mem "Row in memory" :: memory
        memory
            region BUFFER
                var slot -> row
            region TUPLE
                object row : Row : id=50, name=Cy
    end

    link q.n3 --> idx.n0 "scan uses index"
    link idx.n0 --> pg.slot0 "leaf points to page"
    link pg.tuple2 --> mem.row "tuple is a row"
\end{triton}
  \caption{An ``SQL engine anatomy'' poster, authored as a single Triton
source: a query plan, a B+tree index, a slotted storage page, and an in-memory
row, composed in one grid and joined by cross-diagram links. This combination ---
composing heterogeneous diagrams and linking nodes across them --- is what Triton
adds beyond Mermaid.}
\end{figure}

\subsection{What Triton renders}

\begin{itemize}
  \item \textbf{Mermaid-compatible kinds} --- flowchart, sequence, state, class,
        ER, gantt, journey, mindmap, gitgraph, pie, xychart, quadrant, radar,
        sankey, kanban, requirement, block, c4, architecture, packet, timeline.
  \item \textbf{Tree family} (Triton-native) --- a decorated \texttt{tree}
        engine plus value-driven front-ends \texttt{plan}, \texttt{avl},
        \texttt{rbtree}, \texttt{btree}, \texttt{radix}, \texttt{segtree},
        \texttt{heap}. These build \emph{correct-by-construction} structures.
  \item \textbf{Struct / memory family} --- \texttt{array}, \texttt{linkedlist},
        \texttt{memory} (regions with cross-region pointers), \texttt{page}
        (a slotted storage page).
  \item \textbf{Topology} --- cost-tiered node graphs with legends and nested
        groups.
  \item \textbf{Poster composition} --- a grid that embeds any of the above as
        cells, with cell spanning and \emph{cross-diagram links} between panels.
\end{itemize}

\subsection{Design properties}

\begin{description}
  \item[Contracts-first.] The \texttt{Scene} is the one rendering contract; the
        \texttt{DiagramModule} is the one diagram contract.
  \item[Per-diagram IR.] Each kind models its own domain; the god-IR is
        explicitly rejected.
  \item[Grammar = semantics, theme = style.] Source decides \emph{what} a
        diagram means; the theme decides \emph{how} it looks. They never mix.
  \item[Deterministic.] The same source plus theme always yields the same SVG,
        byte for byte (\S\ref{sec:themes-determinism}).
  \item[Correct-by-construction.] The value-driven builders (e.g.\ \texttt{avl})
        compute valid structures from operations; you cannot author an invalid
        one.
\end{description}

\subsection{What Triton is not}

Triton is a text-to-SVG compiler, nothing more. It does \emph{not} ingest data
sources, does \emph{not} construct diagrams from natural-language prompts, does
\emph{not} export PPTX, and has no general animation pipeline. The full
out-of-scope list is in \S\ref{sec:status}.
